% Created 2026-02-24 Tue 22:22
% Intended LaTeX compiler: pdflatex
\documentclass[11pt]{article}
\usepackage[utf8]{inputenc}
\usepackage[T1]{fontenc}
\usepackage{graphicx}
\usepackage{longtable}
\usepackage{wrapfig}
\usepackage{rotating}
\usepackage[normalem]{ulem}
\usepackage{amsmath}
\usepackage{amssymb}
\usepackage{capt-of}
\usepackage{hyperref}
\author{Prologos Project}
\date{2026-02-24}
\title{Logic Engine Design — Comprehensive Implementation Guide\\\medskip
\large From Lattice Foundations to Relational Surface Syntax}
\hypersetup{
 pdfauthor={Prologos Project},
 pdftitle={Logic Engine Design — Comprehensive Implementation Guide},
 pdfkeywords={},
 pdfsubject={},
 pdfcreator={Emacs 30.2 (Org mode 9.7.39)}, 
 pdflang={English}}
\begin{document}

\maketitle
\tableofcontents

\section{Executive Summary}
\label{sec:org94bcbf2}

This document is the implementation guide for Prologos's logic engine — the
substrate on which the relational language (\texttt{defr}, \texttt{rel}, \texttt{solve}) will run.
It follows the research-design-implementation methodology that produced the
Extended Spec Design (\url{2026-02-22\_EXTENDED\_SPEC\_DESIGN.org}).

The design is grounded in two research documents:
\begin{itemize}
\item \href{2026-02-24\_TOWARDS\_A\_GENERAL\_LOGIC\_ENGINE\_ON\_PROPAGATORS.org}{Towards a General Logic Engine on Propagators} — three-layer architecture
\item \href{../research/IMPLEMENTATION\_GUIDE\_CORE\_DS\_PROLOGOS.md}{Implementation Guide: Core Data Structures} — Section 8 (lattice integration)
\end{itemize}

And the relational language vision:
\begin{itemize}
\item \href{principles/RELATIONAL\_LANGUAGE\_VISION.org}{Relational Language Vision} — surface syntax decisions
\end{itemize}

The engine is built bottom-up in 7 phases:
\begin{enumerate}
\item \textbf{Lattice Trait + champ-insert-join} — the algebraic foundation
\item \textbf{Persistent PropNetwork} (Racket-level) — the monotonic data plane as value
\item \textbf{PropNetwork as Prologos Type} — expose network ops to the type system
\item \textbf{UnionFind} — persistent disjoint sets for unification
\item \textbf{Persistent ATMS} — hypothetical reasoning as value
\item \textbf{Tabling} — SLG-style memoization for completeness
\item \textbf{Surface Syntax} — \texttt{defr}, \texttt{rel}, \texttt{solve}, \texttt{\&>}
\end{enumerate}

Key architectural decision: the entire propagator network and ATMS are
\textbf{persistent/immutable values} backed by CHAMP maps. Backtracking = keep
old reference (O(1)). Snapshots = free. Network mobility = serialize value.
LVars are subsumed by PropNetwork cells (join-on-write semantics).

Phases 1-4 are infrastructure with no surface syntax changes. Phase 5-6
add runtime logic capabilities. Phase 7 adds the user-facing language.
\subsection{Infrastructure Gap Analysis}
\label{sec:org4f7517b}

\begin{center}
\begin{tabular}{llll}
Component & Status & What Exists & What's Needed\\
\hline
Persistent collections & COMPLETE & PVec, Map, Set, List, LSeq & —\\
CHAMP lattice helper & COMPLETE & \texttt{champ-insert-join} in champ.rkt & —\\
Transient builders & COMPLETE & TVec, TMap, TSet, with-transient & —\\
Trait system & COMPLETE & Registry, resolution, bundles & Lattice trait + instances\\
Property system & Phase 1 COMPLETE & Storage, flattening, accessors & QuickCheck for :holds (Phase 2+)\\
Spec metadata & Phase 1 COMPLETE & :examples, :deprecated, :doc & :tabled, :answer-mode, :strategy\\
Metavar system & COMPLETE & 4 parallel stores, constraints & Refactor to propagator cells (later)\\
Lattice trait & NOT STARTED & — & trait + 5-8 standard instances\\
Persistent PropNetwork & NOT STARTED & — & Immutable network, pure ops, CHAMP\\
PropNetwork Prologos type & NOT STARTED & — & AST nodes, type rules, reduction\\
UnionFind & NOT STARTED & — & Persistent disjoint sets\\
Persistent ATMS & NOT STARTED & — & Immutable ATMS, worldview switching\\
Tabling & NOT STARTED & — & Producer/consumer, table cells\\
Stratification & NOT STARTED & — & SCC decomposition, stratum evaluation\\
Logic syntax & NOT STARTED & — & defr, rel, solve, \&>, ?var\\
\end{tabular}
\end{center}
\subsection{Critical Path}
\label{sec:org2cd6507}

\begin{verbatim}
Phase 1 (Lattice + champ-insert-join)
  ↓
Phase 2 (Persistent PropNetwork, Racket-level) ←── depends on Phase 1
  ↓
Phase 3 (PropNetwork as Prologos Type) ←── depends on Phase 2
  ↓                    ↓
Phase 4 (UnionFind) ←─ can proceed in parallel with Phase 3
  ↓                    ↓
Phase 5 (Persistent ATMS) ←── depends on Phase 3
  ↓
Phase 6 (Tabling) ←── depends on Phase 3, Phase 5
  ↓
Phase 7 (Surface Syntax) ←── depends on ALL previous phases
\end{verbatim}
\section{Phase 1: Lattice Trait + Standard Instances}
\label{sec:org282f083}

\subsection{1.1 Goal}
\label{sec:orge5d05d7}

Establish the \texttt{Lattice} trait — the algebraic foundation for monotonic
computation. Every propagator cell, LVar, and ATMS label set requires
a lattice domain. By defining \texttt{Lattice} as a standard Prologos trait,
we get automatic dictionary resolution at both compile-time and runtime.
\subsection{1.2 The \texttt{Lattice} Trait}
\label{sec:org2cd22cf}

\begin{verbatim}
;; in lib/prologos/core/lattice.prologos
(ns prologos::core::lattice :no-prelude)

trait Lattice {A : Type}
  bot  : A                     ;; bottom element (no information)
  join : A -> A -> A           ;; least upper bound (merge information)
  leq  : A -> A -> Bool        ;; lattice ordering

;; Laws (to be checked via property testing in Phase 2+):
;;   join is commutative:  join a b = join b a
;;   join is associative:  join (join a b) c = join a (join b c)
;;   join is idempotent:   join a a = a
;;   bot is identity:      join bot a = a
;;   leq is consistent:    leq a b = (join a b == b)
\end{verbatim}
\subsection{1.3 Standard Lattice Instances}
\label{sec:orga808849}

\subsubsection{1.3.1 \texttt{FlatLattice A} — Three-Point Lattice}
\label{sec:org5f53f1b}

The simplest useful lattice: bottom → exactly one value → top (contradiction).
Used for basic unification cells where a variable is either unknown, bound to
exactly one value, or in conflict.

\begin{verbatim}
;; in lib/prologos/core/lattice-instances.prologos

;; FlatLattice wraps any Eq type into a flat lattice
;; ⊥ → value → ⊤
functor FlatLattice {A : Type}
  :where (Eq A)
  :unfolds <(Option [Either A Unit])>
  ;; none = ⊥, some (left v) = value, some (right unit) = ⊤

instance Lattice [FlatLattice A] where (Eq A)
  bot  = none
  join = flat-join
  leq  = flat-leq

defn flat-join [a b]
  match [a b]
    | [none _]                     -> b
    | [_ none]                     -> a
    | [[some [left v1]] [some [left v2]]]
      -> if [eq? v1 v2] a [some [right unit]]
    | _                            -> [some [right unit]]
\end{verbatim}
\subsubsection{1.3.2 \texttt{SetLattice A} — Powerset Lattice (Set Union)}
\label{sec:orge1c4088}

Growing sets: elements can only be added, never removed. Join = union.
Used for Herbrand interpretations (sets of ground facts) and LVar-Sets.

\begin{verbatim}
;; SetLattice: growing sets with union as join
functor SetLattice {A : Type}
  :unfolds [Set A]

instance Lattice [SetLattice A]
  bot  = #{}
  join = set-union
  leq  = set-subset?
\end{verbatim}
\subsubsection{1.3.3 \texttt{MapLattice K V} — Pointwise Map Lattice}
\label{sec:org801d66f}

Maps where each key has an independent lattice-valued value. Join =
pointwise join of values. Used for LVar-Maps and substitution stores.

\begin{verbatim}
;; MapLattice: pointwise lattice on maps
functor MapLattice {K V : Type}
  :where (Lattice V)
  :unfolds [Map K V]

instance Lattice [MapLattice K V] where (Lattice V)
  bot  = {}
  join = map-lattice-join
  leq  = map-lattice-leq

defn map-lattice-join [m1 m2]
  ;; For keys in both: join values. For keys in one: include.
  [merge-with join m1 m2]
\end{verbatim}
\subsubsection{1.3.4 \texttt{IntervalLattice} — Numeric Intervals}
\label{sec:org71864ba}

Intervals \texttt{[lo, hi]} with intersection as meet, hull as join. Used for
numeric constraint propagation (CLP(R/FD)).

\begin{verbatim}
;; IntervalLattice for numeric constraints
deftype Interval
  mk-interval : Rat -> Rat -> Interval    ;; lo, hi (lo <= hi)
  ;; Special: bot = (-∞, +∞) = unconstrained
  ;; Special: top = empty = contradicted

instance Lattice Interval
  bot  = interval-unconstrained
  join = interval-intersect    ;; narrowing: more constrained = higher
  leq  = interval-subsumes?   ;; [a,b] ⊑ [c,d] iff [c,d] ⊆ [a,b]
\end{verbatim}
\subsubsection{1.3.5 \texttt{BoolLattice} — Two-Point Lattice}
\label{sec:orga024ff8}

The simplest non-trivial lattice: false < true. Join = OR.
Used as building block and for boolean constraint propagation.

\begin{verbatim}
instance Lattice Bool
  bot  = false
  join = bool-or
  leq  = bool-implies?
\end{verbatim}
\subsection{1.4 Racket-Level Implementation}
\label{sec:org881815e}

At the Racket level, lattice operations are dispatched via the existing trait
system. The \texttt{Lattice} trait follows the standard pattern:

\begin{center}
\begin{tabular}{ll}
Component & How\\
\hline
Trait struct & \texttt{trait-meta} in \texttt{current-trait-registry}\\
Instances & \texttt{(impl-entry 'Lattice ...)} in \texttt{current-impl-registry}\\
Dictionary type & Sigma type: \texttt{(bot . (join . leq))} (3 methods)\\
Resolution & Standard \texttt{resolve-trait-constraints!} pipeline\\
\end{tabular}
\end{center}

No new AST nodes needed for Phase 1. The trait and instance declarations
use existing infrastructure.
\subsection{1.5 \texttt{champ-insert-join} — Racket-Level Helper}
\label{sec:org4e9ade6}

A lattice-aware CHAMP insert that joins on collision, rather than replacing.
This is the foundational Racket-level helper that enables the PropNetwork's
cells map to do lattice merge natively.

\begin{verbatim}
;; In racket/prologos/champ.rkt (already implemented)

;; champ-insert-join : champ-root hash key val (val val → val) → champ-root
(define (champ-insert-join root hash key val join-fn)
  (let ([existing (champ-lookup root hash key)])
    (if (eq? existing 'none)
        (champ-insert root hash key val)
        (champ-insert root hash key (join-fn existing val)))))
\end{verbatim}

This enables the propagator network to store cell values in a CHAMP map
with join-on-write semantics: writing to a cell computes
\texttt{merge-fn(old-value, new-value)} and inserts the result.
\subsection{1.6 New Files}
\label{sec:orgfdef308}

\begin{center}
\begin{tabular}{ll}
File & Purpose\\
\hline
\texttt{lib/prologos/core/lattice.prologos} & \texttt{Lattice} trait declaration\\
\texttt{lib/prologos/core/lattice-instances.prologos} & Standard instances\\
\texttt{tests/test-lattice.rkt} & Trait resolution + law tests\\
\end{tabular}
\end{center}
\subsection{1.7 Tests (\textasciitilde{}25)}
\label{sec:orgd3cc659}

\begin{itemize}
\item Lattice trait registered correctly
\item FlatLattice: bot is identity, join of same values, join of different → top
\item SetLattice: bot = empty, join = union, leq = subset
\item MapLattice: pointwise join, partial maps merge correctly
\item BoolLattice: basic operations
\item Trait resolution: \texttt{(Lattice [FlatLattice Nat])} resolves correctly
\item Laws: commutativity, associativity, idempotency of join (6 tests per instance)
\end{itemize}
\subsection{1.8 Dependencies}
\label{sec:orgac4a5a8}

\begin{itemize}
\item Existing trait system (\texttt{process-trait}, \texttt{process-impl})
\item Existing collection types (\texttt{Set}, \texttt{Map} for SetLattice, MapLattice)
\item \texttt{Eq} trait (for \texttt{FlatLattice})
\item \texttt{champ-insert-join} in \texttt{champ.rkt} (already implemented)
\item \textbf{No new AST nodes}
\item \textbf{No changes to typing-core.rkt}
\end{itemize}

---
\section{Phase 2: Persistent Propagator Network}
\label{sec:org440226c}

\subsection{2.1 Goal}
\label{sec:org524bef5}

Implement the monotonic data plane as a \textbf{persistent, immutable value}.
The entire propagator network — cells, propagators, worklist, and
metadata — is a single Racket struct backed by CHAMP maps. All operations
are pure functions: they take a network and return a new network.

This is a critical design choice: the propagator network is a \textbf{first-class
value} that can be snapshotted (free), backtracked (O(1) — keep old
reference), migrated (serialize and send), and compared (structural equality).
\subsection{2.2 Architecture}
\label{sec:orgf342acb}

\begin{verbatim}
┌───────────────────────────────────────┐
│  prop-network (persistent value)      │
│                                       │
│  cells : champ-root (cell-id → cell)  │
│  propagators : champ-root (prop-id)   │
│  worklist : list of prop-id           │
│  merge-fns : champ-root (cell-id → fn)│
│  fuel : Nat                           │
│  contradiction : #f | cell-id         │
│  next-cell-id : Nat                   │
│  next-prop-id : Nat                   │
│                                       │
│  All operations return NEW networks:  │
│  net-new-cell   : Net → ... → (Net,Id)│
│  net-cell-write : Net → Id → V → Net  │
│  run-to-quiescence : Net → Net        │
└───────────────────────────────────────┘
\end{verbatim}

Backtracking model:
\begin{verbatim}
let net0 = (make-prop-network)
let net1 = (net-cell-write net0 cell-a 42)
let net2 = (net-cell-write net1 cell-b 99)

;; Backtrack to net1: simply use the old reference
;; net0 and net1 are unchanged — structural sharing via CHAMP
\end{verbatim}
\subsection{2.3 Core Data Structures (All Persistent)}
\label{sec:org90e356a}

\subsubsection{2.3.1 Identity Types}
\label{sec:org7a4255b}

\begin{verbatim}
;; In racket/prologos/propagator.rkt

;; Deterministic counters (no gensym — no global state)
(struct cell-id (n) #:transparent)
(struct prop-id (n) #:transparent)
\end{verbatim}

Cell and propagator identities are monotonic counters \emph{inside} the network.
This makes networks deterministic (no gensym side effects) and serializable.
\subsubsection{2.3.2 \texttt{prop-cell} — Propagator Cell (Immutable)}
\label{sec:orgac05076}

\begin{verbatim}
(struct prop-cell
  (value        ;; Expr — current lattice value (starts at bot)
   dependents)  ;; champ-root (set of prop-id → #t)
  #:transparent)
\end{verbatim}

Note: no \texttt{id} field in the cell struct itself — the identity is the key
in the network's cells map. No \texttt{domain} field — the merge function is
stored in the network's \texttt{merge-fns} map, keyed by cell-id.
\subsubsection{2.3.3 \texttt{propagator} — Monotone Function (Immutable)}
\label{sec:org249cb9e}

\begin{verbatim}
(struct propagator
  (inputs      ;; list of cell-id
   outputs     ;; list of cell-id
   fire-fn)    ;; (prop-network → prop-network) — pure state transformer
  #:transparent)
\end{verbatim}

The \texttt{fire-fn} is a \textbf{pure function} from network to network. It reads input
cells from the network, computes new values, and returns a network with
updated output cells. No side effects.
\subsubsection{2.3.4 \texttt{prop-network} — The Network as Value}
\label{sec:orga8600cc}

\begin{verbatim}
(struct prop-network
  (cells            ;; champ-root : cell-id → prop-cell
   propagators      ;; champ-root : prop-id → propagator
   worklist         ;; list of prop-id (ephemeral; empty at quiescence)
   next-cell-id     ;; Nat — monotonic counter for cell ids
   next-prop-id     ;; Nat — monotonic counter for propagator ids
   fuel             ;; Nat — step limit to prevent runaway
   contradiction    ;; #f | cell-id — first contradiction encountered
   merge-fns        ;; champ-root : cell-id → (Expr Expr → Expr)
   contradiction-fns) ;; champ-root : cell-id → (Expr → Bool), optional
  #:transparent)
\end{verbatim}

Key design choices:
\begin{itemize}
\item \textbf{CellId = Nat counter}: deterministic, no gensym, serializable
\item \textbf{merge-fns} stored per-cell: captures the lattice's \texttt{join} at cell creation
\item \textbf{contradiction-fns} stored per-cell: optional predicate that detects if a cell
value represents contradiction (\texttt{top}). This avoids requiring a \texttt{top?} method
in the \texttt{Lattice} trait or introducing a \texttt{BoundedLattice} sub-trait (Prologos
has \textbf{no trait hierarchies} — see DESIGN\textsubscript{PRINCIPLES.org}). Instead, contradiction
detection is a per-cell concern supplied at cell creation time.
\item \textbf{worklist} is a plain Racket list: ephemeral within \texttt{run-to-quiescence},
always empty in saved snapshots (quiescent networks have empty worklists)
\item \textbf{All CHAMP-backed}: O(log₃₂ n) ≈ O(C) where C ≤ 7 for practical n
\end{itemize}
\subsection{2.4 Pure Operations}
\label{sec:orgfce1fa2}

All operations take a network and return a new network. The old network
is never modified.

\begin{verbatim}
;; Create empty network
(define (make-prop-network [fuel 1000000])
  (prop-network champ-empty    ;; cells
                champ-empty    ;; propagators
                '()            ;; worklist
                0              ;; next-cell-id
                0              ;; next-prop-id
                fuel           ;; fuel
                #f             ;; no contradiction
                champ-empty    ;; merge-fns
                champ-empty))  ;; contradiction-fns

;; Add a new cell: returns (values new-network cell-id)
;; contradicts? is optional (#f = no contradiction detection at this cell)
(define (net-new-cell net initial-value merge-fn [contradicts? #f])
  (define id (cell-id (prop-network-next-cell-id net)))
  (define cell (prop-cell initial-value champ-empty))
  (define net*
    (struct-copy prop-network net
      [cells (champ-insert (prop-network-cells net)
                            (equal-hash-code id) id cell)]
      [merge-fns (champ-insert (prop-network-merge-fns net)
                                (equal-hash-code id) id merge-fn)]
      [next-cell-id (+ 1 (prop-network-next-cell-id net))]))
  (values
   (if contradicts?
       (struct-copy prop-network net*
         [contradiction-fns
          (champ-insert (prop-network-contradiction-fns net*)
                        (equal-hash-code id) id contradicts?)])
       net*)
   id))

;; Read a cell's value
(define (net-cell-read net cid)
  (define cell (champ-lookup (prop-network-cells net)
                              (equal-hash-code cid) cid))
  (if (eq? cell 'none)
      (error 'net-cell-read "unknown cell: ~a" cid)
      (prop-cell-value cell)))

;; Write a cell: computes merge-fn(old, new), enqueues dependents if changed
(define (net-cell-write net cid new-val)
  (define cells (prop-network-cells net))
  (define cell (champ-lookup cells (equal-hash-code cid) cid))
  (when (eq? cell 'none)
    (error 'net-cell-write "unknown cell: ~a" cid))
  (define merge-fn (champ-lookup (prop-network-merge-fns net)
                                  (equal-hash-code cid) cid))
  (define old-val (prop-cell-value cell))
  (define merged (merge-fn old-val new-val))
  (if (equal? merged old-val)
      net  ;; No change — return same network
      (let* ([new-cell (struct-copy prop-cell cell [value merged])]
             [new-cells (champ-insert cells (equal-hash-code cid) cid new-cell)]
             ;; Enqueue dependents
             [deps (champ-keys (prop-cell-dependents cell))]
             [new-wl (append deps (prop-network-worklist net))])
        (struct-copy prop-network net
          [cells new-cells]
          [worklist new-wl]))))

;; Add a propagator: returns (values new-network prop-id)
(define (net-add-propagator net input-ids output-ids fire-fn)
  (define pid (prop-id (prop-network-next-prop-id net)))
  (define prop (propagator input-ids output-ids fire-fn))
  ;; Register pid as dependent of each input cell
  (define new-cells
    (for/fold ([cells (prop-network-cells net)])
              ([cid (in-list input-ids)])
      (define cell (champ-lookup cells (equal-hash-code cid) cid))
      (if (eq? cell 'none) cells
          (let ([new-deps (champ-insert (prop-cell-dependents cell)
                                         (equal-hash-code pid) pid #t)])
            (champ-insert cells (equal-hash-code cid) cid
                          (struct-copy prop-cell cell [dependents new-deps]))))))
  (values
   (struct-copy prop-network net
     [cells new-cells]
     [propagators (champ-insert (prop-network-propagators net)
                                 (equal-hash-code pid) pid prop)]
     [next-prop-id (+ 1 (prop-network-next-prop-id net))]
     ;; Schedule initial firing
     [worklist (cons pid (prop-network-worklist net))])
   pid))
\end{verbatim}
\subsection{2.5 Concrete \texttt{fire-fn} Example: Adder Propagator}
\label{sec:org437d837}

To ground the abstract \texttt{fire-fn} concept, here is a concrete example: a
propagator that adds two cells and writes the result to a third.

\begin{verbatim}
;; adder-fire-fn : cell-id cell-id cell-id → (prop-network → prop-network)
;; Reads cells A and B, writes A+B to cell C
(define (make-adder-fire-fn a-id b-id c-id)
  (lambda (net)
    (define a-val (net-cell-read net a-id))
    (define b-val (net-cell-read net b-id))
    ;; Only propagate if both inputs are known (not bot)
    (if (and (not (eq? a-val 'bot)) (not (eq? b-val 'bot)))
        (net-cell-write net c-id (+ a-val b-val))
        net)))  ;; No change — return same network

;; Usage: wire up an adder network
(define net0 (make-prop-network))
(define-values (net1 cell-a) (net-new-cell net0 'bot flat-join))
(define-values (net2 cell-b) (net-new-cell net1 'bot flat-join))
(define-values (net3 cell-c) (net-new-cell net2 'bot flat-join))
(define-values (net4 adder-pid)
  (net-add-propagator net3
    (list cell-a cell-b)   ;; inputs
    (list cell-c)          ;; outputs
    (make-adder-fire-fn cell-a cell-b cell-c)))

;; Write to inputs, run to quiescence
(define net5 (net-cell-write net4 cell-a 3))
(define net6 (net-cell-write net5 cell-b 4))
(define net7 (run-to-quiescence net6))
(net-cell-read net7 cell-c) ;; → 7
\end{verbatim}

Key observations:
\begin{itemize}
\item The \texttt{fire-fn} is a \textbf{closure} capturing cell-ids, returning a \textbf{pure function}
from network to network
\item It reads inputs via \texttt{net-cell-read} (from the network argument, not side effects)
\item It writes via \texttt{net-cell-write} (returns new network, not mutation)
\item If inputs are insufficient (\texttt{bot}), it returns the network unchanged (idempotent)
\item The same \texttt{fire-fn} works correctly regardless of scheduling order
\end{itemize}

For the logic engine (Phase 7), a typical \texttt{fire-fn} will be a \emph{clause propagator}
that attempts unification and writes answer substitutions to a table cell. See
Appendix A: Resolution by Example for a worked walkthrough.
\subsection{2.6 \texttt{run-to-quiescence} — Pure Loop}
\label{sec:orgcd56e78}

The scheduler is not a separate struct — it is the \texttt{worklist} field of the
network itself. \texttt{run-to-quiescence} is a pure tail-recursive function.

\begin{verbatim}
(define (run-to-quiescence net)
  (cond
    ;; Already contradicted — stop
    [(prop-network-contradiction net) net]
    ;; Fuel exhausted — stop
    [(<= (prop-network-fuel net) 0) net]
    ;; Worklist empty — quiescent (fixed point reached)
    [(null? (prop-network-worklist net)) net]
    ;; Fire next propagator
    [else
     (let* ([pid (car (prop-network-worklist net))]
            [rest (cdr (prop-network-worklist net))]
            [net* (struct-copy prop-network net
                    [worklist rest]
                    [fuel (sub1 (prop-network-fuel net))])]
            [prop (champ-lookup (prop-network-propagators net*)
                                (equal-hash-code pid) pid)])
       (if (eq? prop 'none)
           ;; Propagator removed or unknown — skip
           (run-to-quiescence net*)
           ;; Fire: pure function from network to network
           (run-to-quiescence ((propagator-fire-fn prop) net*))))]))
\end{verbatim}

Properties:
\begin{itemize}
\item \textbf{Pure}: input is a \texttt{prop-network}, output is a \texttt{prop-network}
\item \textbf{Deterministic}: same inputs → same fixed point (lattice commutativity)
\item \textbf{Convergent}: finite-height lattices guarantee termination
\item \textbf{Fuel-limited}: prevents runaway; configurable via \texttt{fuel} field
\item \textbf{Backtrackable}: keep old \texttt{net} reference = instant O(1) backtrack
\end{itemize}
\subsection{2.7 Contradiction Handling (Per-Cell \texttt{contradicts?} Predicate)}
\label{sec:orgf642aaa}

Contradiction is detected via a \textbf{per-cell predicate}, not a trait method.
Each cell may optionally have a \texttt{contradicts?} function (supplied at creation
time via \texttt{net-new-cell}) that checks if a value represents inconsistency.

This design avoids requiring \texttt{top?} in the \texttt{Lattice} trait (which would force
every lattice to define it, even those without a sensible top element) and
avoids introducing a \texttt{BoundedLattice} sub-trait (Prologos has no trait
hierarchies — see DESIGN\textsubscript{PRINCIPLES.org} § "No Trait Hierarchies").

\begin{verbatim}
;; In net-cell-write, after computing merged:
(define (net-cell-write net cid new-val)
  ;; ... (merge as before) ...
  ;; After merge, check for contradiction:
  (define cfn (champ-lookup (prop-network-contradiction-fns net)
                             (equal-hash-code cid) cid))
  (define contradicted?
    (and (not (eq? cfn 'none))   ;; cell has a contradicts? fn
         (cfn merged)))          ;; the merged value is contradictory
  (if contradicted?
      (struct-copy prop-network net* [contradiction cid])
      net*))

;; Example contradiction predicates for standard lattices:
;; FlatLattice: (lambda (v) (and (some? v) (right? (unwrap v))))  ;; some(right(unit)) = top
;; IntervalLattice: (lambda (v) (interval-empty? v))              ;; empty interval = contradicted
;; SetLattice: #f  ;; no contradiction (sets only grow, never contradict)
\end{verbatim}

At the pure-network level, contradiction simply records which cell was
contradicted. The ATMS layer (Phase 5) interprets this and performs
dependency-directed backtracking. Without ATMS, contradiction halts the
network.
\subsection{2.8 LVars Are Subsumed by Cells}
\label{sec:orgb8e4f7d}

A key simplification: \textbf{LVars are just cells in a persistent network}.

\begin{center}
\begin{tabular}{ll}
LVar operation & Network equivalent\\
\hline
\texttt{lvar-new} & \texttt{net-new-cell} with appropriate \texttt{merge-fn}\\
\texttt{lvar-put} & \texttt{net-cell-write} (already does join-on-write)\\
\texttt{lvar-get} & \texttt{net-cell-read} (synchronous in single-thread)\\
\texttt{lvar-freeze} & Read cell value from a quiescent network\\
\texttt{lvar-set-add} & \texttt{net-cell-write} on cell with \texttt{SetLattice}\\
\texttt{lvar-map-put} & \texttt{net-cell-write} on cell with \texttt{MapLattice}\\
\end{tabular}
\end{center}

This elimination removes the need for a separate LVar module and \textasciitilde{}10
dedicated LVar AST nodes. The propagator network \emph{is} the lattice-compatible
collection — its cells natively do join-on-write via the \texttt{merge-fn}.

Threshold reads: for the synchronous logic engine (not parallel), threshold
reads are implemented as propagators that check conditions after quiescence.
True blocking threshold reads (for parallel LVar-style programming) are
deferred to when actor/place integration is built.
\subsection{2.9 AST Nodes (\textasciitilde{}12)}
\label{sec:org958619c}

New \texttt{expr-*} structs in \texttt{syntax.rkt}:

\begin{center}
\begin{tabular}{lll}
Node & Fields & Semantics\\
\hline
\texttt{expr-prop-network} & net-value & Runtime wrapper for network\\
\texttt{expr-cell-id} & n & Cell identity (Nat counter)\\
\texttt{expr-net-new} & fuel-expr & Create empty network\\
\texttt{expr-net-new-cell} & net, init-val, merge & Add cell, return (net, id)\\
\texttt{expr-net-cell-read} & net, cell-id & Read cell value\\
\texttt{expr-net-cell-write} & net, cell-id, val & Write (merge) value\\
\texttt{expr-net-add-prop} & net, ins, outs, fn & Add propagator\\
\texttt{expr-net-run} & net & Run to quiescence\\
\texttt{expr-net-snapshot} & net & Snapshot (identity — free)\\
\texttt{expr-net-contradict?} & net & Check for contradiction\\
\texttt{expr-net-type} & — & Type constructor PropNetwork\\
\texttt{expr-cell-id-type} & — & Type constructor CellId\\
\end{tabular}
\end{center}
\subsection{2.10 New/Modified Files}
\label{sec:orga84ceb3}

\begin{center}
\begin{tabular}{ll}
File & Changes\\
\hline
\texttt{racket/prologos/propagator.rkt} & NEW: Persistent network, pure ops\\
\texttt{racket/prologos/syntax.rkt} & +12 AST nodes\\
\texttt{racket/prologos/typing-core.rkt} & Infer/check rules for network ops\\
\texttt{racket/prologos/reduction.rkt} & Reduction rules for network ops\\
\texttt{racket/prologos/elaborator.rkt} & Surface → core for network ops\\
\texttt{racket/prologos/zonk.rkt} & Zonk cases for new nodes\\
\texttt{racket/prologos/pretty-print.rkt} & PP cases for new nodes\\
\texttt{racket/prologos/qtt.rkt} & Multiplicity for network ops\\
\texttt{racket/prologos/substitution.rkt} & Subst cases for new nodes\\
\texttt{racket/prologos/unify.rkt} & Unify cases for new nodes\\
\texttt{tests/test-propagator.rkt} & NEW\\
\texttt{tests/test-propagator-network.rkt} & NEW\\
\texttt{tests/test-propagator-quiescence.rkt} & NEW\\
\end{tabular}
\end{center}
\subsection{2.11 Tests (\textasciitilde{}60)}
\label{sec:org35038c8}

\begin{itemize}
\item Network creation: empty, with fuel
\item Cell creation: returns new network + cell-id
\item Cell read: returns value
\item Cell write: merge produces join (FlatLattice, SetLattice)
\item Cell write: no change when value doesn't increase → same network returned
\item Cell write: enqueues dependents
\item Propagator: two cells connected, write propagates
\item Propagator: diamond network (A → B, A → C, B+C → D)
\item run-to-quiescence: fires to fixed point
\item run-to-quiescence: fuel limit prevents runaway
\item run-to-quiescence: contradiction recorded in network
\item Persistence: old network unchanged after write
\item Persistence: old network unchanged after run-to-quiescence
\item Snapshot: identity operation (\texttt{net} is already persistent)
\item Backtracking: use old reference after contradiction
\item LVar-style: cell with SetLattice merge, add elements, freeze via read
\item LVar-style: cell with MapLattice merge, pointwise join
\item Integration: lattice instances work through network cells
\end{itemize}
\subsection{2.12 Dependencies}
\label{sec:org9aa4f6d}

\begin{itemize}
\item Phase 1 (Lattice trait and instances)
\item \texttt{champ-insert-join} in \texttt{champ.rkt} (used internally for cell merge)
\item Existing trait resolution for Lattice dictionary dispatch
\item 12+ files modified (standard AST node pipeline)
\end{itemize}

---
\section{Phase 3: PropNetwork as Prologos Type}
\label{sec:org2417eda}

\subsection{3.1 Goal}
\label{sec:orgd4e691f}

Expose the Racket-level persistent PropNetwork to Prologos's type system.
This phase adds the 12 AST nodes defined in Phase 2's design, threading
them through the full 12-file pipeline (syntax → typing-core → reduction →
elaborator → zonk → pretty-print → qtt → substitution → unify → tests).
\subsection{3.2 Why a Separate Phase}
\label{sec:org538889b}

Phase 2 implements the Racket-level data structures and algorithms. Phase 3
wires them into Prologos as first-class values. This separation follows
the established pattern: Racket infrastructure first, then Prologos type
system integration.
\subsection{3.3 Type Signatures}
\label{sec:org9270790}

\begin{verbatim}
;; PropNetwork is an opaque type wrapping the persistent network
deftype PropNetwork
  mk-prop-network : PropNetwork

;; CellId identifies a cell within a network
deftype CellId
  mk-cell-id : Nat -> CellId

;; Operations as typed functions
spec net-new : Nat -> PropNetwork
spec net-new-cell : {A : Type} where (Lattice A)
  PropNetwork -> A -> (A -> A -> A) -> [PropNetwork * CellId]
spec net-cell-read : {A : Type} PropNetwork -> CellId -> A
spec net-cell-write : {A : Type} PropNetwork -> CellId -> A -> PropNetwork
spec net-add-propagator : PropNetwork -> [List CellId] -> [List CellId]
  -> [PropNetwork -> PropNetwork] -> [PropNetwork * PropId]
spec net-run : PropNetwork -> PropNetwork
spec net-snapshot : PropNetwork -> PropNetwork
spec net-contradict? : PropNetwork -> Bool
\end{verbatim}
\subsection{3.4 LVar Operations as Library Functions}
\label{sec:org9b3cf28}

LVar-style operations are library functions on top of PropNetwork, not
separate AST nodes:

\begin{verbatim}
;; A "growing set" = a cell with SetLattice merge
spec make-set-cell : {A : Type} where (Eq A)
  PropNetwork -> [PropNetwork * CellId]
defn make-set-cell [net]
  [net-new-cell net #{} set-union]

;; A "growing map" = a cell with MapLattice merge
spec make-map-cell : {K V : Type} where (Eq K) (Lattice V)
  PropNetwork -> [PropNetwork * CellId]
defn make-map-cell [net]
  [net-new-cell net {} [merge-with join]]
\end{verbatim}
\subsection{3.5 AST Nodes (12)}
\label{sec:org298d62d}

Same 12 nodes as Phase 2 design (\texttt{expr-prop-network}, \texttt{expr-cell-id},
\texttt{expr-net-new}, \texttt{expr-net-new-cell}, \texttt{expr-net-cell-read},
\texttt{expr-net-cell-write}, \texttt{expr-net-add-prop}, \texttt{expr-net-run},
\texttt{expr-net-snapshot}, \texttt{expr-net-contradict?}, \texttt{expr-net-type},
\texttt{expr-cell-id-type}).
\subsection{3.6 New/Modified Files}
\label{sec:org3672282}

\begin{center}
\begin{tabular}{ll}
File & Changes\\
\hline
\texttt{racket/prologos/syntax.rkt} & +12 AST nodes\\
\texttt{racket/prologos/typing-core.rkt} & Infer/check rules for network ops\\
\texttt{racket/prologos/reduction.rkt} & Reduction rules (delegate to .rkt)\\
\texttt{racket/prologos/elaborator.rkt} & Surface → core translation\\
\texttt{racket/prologos/zonk.rkt} & Zonk cases\\
\texttt{racket/prologos/pretty-print.rkt} & PP cases\\
\texttt{racket/prologos/qtt.rkt} & Multiplicity (:w for networks)\\
\texttt{racket/prologos/substitution.rkt} & Subst cases\\
\texttt{racket/prologos/unify.rkt} & Unify cases\\
\texttt{lib/prologos/core/propagator.prologos} & NEW: Library wrappers (set/map cells)\\
\texttt{tests/test-propagator-types.rkt} & NEW\\
\texttt{tests/test-propagator-integration.rkt} & NEW\\
\texttt{tests/test-propagator-lvar.rkt} & NEW: LVar-style cell tests\\
\end{tabular}
\end{center}
\subsection{3.7 Tests (\textasciitilde{}50)}
\label{sec:org370b5b5}

\begin{itemize}
\item Type checking: \texttt{net-new} returns \texttt{PropNetwork}
\item Type checking: \texttt{net-cell-read} returns correct element type
\item Type checking: \texttt{net-cell-write} returns \texttt{PropNetwork}
\item Eval: create network, add cells, read values
\item Eval: write value, verify merge
\item Eval: propagator fires on cell write
\item Eval: diamond network (A → B, A → C, B+C → D)
\item Eval: run-to-quiescence reaches fixed point
\item Eval: fuel exhaustion
\item Eval: contradiction detection
\item Persistence: old network unchanged after operations
\item LVar-style: set cell grows via union
\item LVar-style: map cell does pointwise join
\item LVar-style: freeze = read from quiescent network
\item Integration: Lattice trait resolution for cell merge functions
\end{itemize}
\subsection{3.8 Dependencies}
\label{sec:org19ed3a6}

\begin{itemize}
\item Phase 2 (Racket-level PropNetwork implementation)
\item Phase 1 (Lattice trait for type-level merge function constraints)
\end{itemize}

---
\section{Phase 4: UnionFind — Persistent Disjoint Sets}
\label{sec:orgcf74b53}

\subsection{4.1 Goal}
\label{sec:orgbe13628}

Implement a persistent union-find data structure (Conchon \& Filliâtre 2007)
with backtracking support. This is the core data structure for unification
in the logic engine. Unlike the current metavar store (mutable hash table),
a persistent union-find supports efficient backtracking for search.
\subsection{4.2 Design}
\label{sec:org8af1162}

The union-find uses path splitting (not path compression) to maintain
persistence. Each node stores a parent reference and a rank. Find follows
parent pointers to the root. Union merges by rank.

\begin{verbatim}
;; In racket/prologos/union-find.rkt

(struct uf-store
  (nodes)     ;; persistent map: id → uf-node
  #:transparent)

(struct uf-node
  (parent     ;; id (self = root)
   rank       ;; natural number
   value)     ;; optional payload (for unification: the term)
  #:transparent)

;; Create empty store
(define (uf-empty) (uf-store (hasheq)))

;; Make a new set with one element
(define (uf-make-set store id value)
  (uf-store
   (hash-set (uf-store-nodes store) id
             (uf-node id 0 value))))

;; Find root of id (with path splitting for amortized cost)
(define (uf-find store id)
  ;; Returns (values root-id updated-store)
  ;; Path splitting: point each node to its grandparent
  ...)

;; Union two sets, merge by rank
(define (uf-union store id1 id2 merge-fn)
  ;; Find roots, merge smaller into larger
  ;; merge-fn combines values if both have payload
  ;; Returns updated store
  ...)

;; Snapshot for backtracking
(define (uf-snapshot store) store)  ;; Already persistent!
\end{verbatim}
\subsection{4.3 Key Properties}
\label{sec:orgc31867b}

\begin{itemize}
\item \textbf{Persistent}: Union/find return new stores, old stores unchanged
\item \textbf{Backtrackable}: Save a reference to the old store = instant backtrack
\item \textbf{Efficient}: O(log n) find and union (path splitting without compression)
\item \textbf{Value-carrying}: Nodes carry optional payloads (unified terms)
\end{itemize}
\subsection{4.4 Integration with Logic Engine: UF vs Cell Division of Labor}
\label{sec:org13df068}

The union-find and propagator cells serve \textbf{complementary} roles in the logic
engine. They are not alternatives — both are needed:

\begin{center}
\begin{tabular}{lll}
Operation & Uses UF & Uses PropNetwork Cell\\
\hline
\texttt{?x = ?y} (var-var) & \texttt{uf-union x y} & —\\
\texttt{?x = 42} (var-value) & \texttt{uf-find x} → set value & —\\
\texttt{?x = f(?y, ?z)} (complex) & \texttt{uf-union} + subterm cells & Cell per subterm for propagation\\
Lattice-valued accumulation & — & \texttt{net-cell-write} with join\\
Table answer sets & — & Cell with \texttt{SetLattice} merge\\
Numeric constraints & — & Cell with \texttt{IntervalLattice}\\
Backtracking & Keep old \texttt{uf-store} & Keep old \texttt{prop-network}\\
\end{tabular}
\end{center}

The division of labor:
\begin{itemize}
\item \textbf{UnionFind}: structural equality — discovering that two variables must refer
to the same term. This is the fast path for standard Prolog-style unification.
O(log n) find-with-path-splitting.
\item \textbf{PropNetwork cells}: lattice-valued accumulation — aggregating information that
can only grow (sets of answers, numeric intervals, constraint domains). This is
the slow-but-general path for constraint propagation.
\end{itemize}

In practice, a logic engine query creates \emph{both}:
\begin{enumerate}
\item A \texttt{uf-store} for the substitution (logic variable bindings)
\item A \texttt{prop-network} for constraint cells (tables, aggregates, numeric domains)
\end{enumerate}

Both are persistent, both support O(1) backtracking. The solver threads both
through computation as a pair: \texttt{(uf-store, prop-network)}.
\subsection{4.5 AST Nodes (\textasciitilde{}6)}
\label{sec:org5753c87}

\begin{center}
\begin{tabular}{lll}
Node & Fields & Semantics\\
\hline
\texttt{expr-uf-empty} & — & Create empty union-find\\
\texttt{expr-uf-make-set} & store, id, val & Add new singleton set\\
\texttt{expr-uf-find} & store, id & Find root of set\\
\texttt{expr-uf-union} & store, id1, id2 & Union two sets\\
\texttt{expr-uf-value} & store, id & Get value at id's root\\
\texttt{expr-uf-type} & — & Type constructor UnionFind\\
\end{tabular}
\end{center}
\subsection{4.6 New/Modified Files}
\label{sec:org71a450b}

\begin{center}
\begin{tabular}{ll}
File & Changes\\
\hline
\texttt{racket/prologos/union-find.rkt} & NEW: Persistent UF\\
\texttt{racket/prologos/syntax.rkt} & +6 AST nodes\\
+ standard pipeline files & Type rules, reduction, etc.\\
\texttt{tests/test-union-find.rkt} & NEW\\
\end{tabular}
\end{center}
\subsection{4.7 Tests (\textasciitilde{}30)}
\label{sec:org50f3ade}

\begin{itemize}
\item Empty store creation
\item Make-set: adds node
\item Find: returns self for singleton
\item Union: merges two sets
\item Union: rank-based merging
\item Find after union: both ids → same root
\item Value: retrieve payload after union
\item Persistence: old store unchanged after union
\item Backtracking: restore snapshot works
\item Performance: 1000 unions, find all roots
\item Integration: as substitution store for simple unification
\end{itemize}
\subsection{4.8 Dependencies}
\label{sec:org0f94316}

\begin{itemize}
\item None (self-contained data structure)
\item Can proceed in parallel with Phase 3
\end{itemize}

---
\section{Phase 5: Persistent ATMS Layer — Hypothetical Reasoning}
\label{sec:org4127204}

\subsection{5.1 Goal}
\label{sec:orgb674077}

Implement the Assumption-Based Truth Maintenance System (ATMS) as a
\textbf{persistent, immutable value}. Like the propagator network, the ATMS is
backed entirely by CHAMP maps. Backtracking = use old reference.
Switching worldviews = \texttt{struct-copy} with new \texttt{believed} set.

This validates the "Multiverse Mechanism" from the propagator research:
choice-point forking maps directly onto ATMS worldview management.
\subsection{5.2 Core Data Structures (All Persistent)}
\label{sec:orgc5fea09}

\subsubsection{\texttt{assumption} — Hypothetical Premise}
\label{sec:orgc1d5fd5}

\begin{verbatim}
;; In racket/prologos/atms.rkt

(struct assumption-id (n) #:transparent)

(struct assumption
  (name       ;; symbol (for display)
   datum)     ;; optional: the value this assumption asserts
  #:transparent)
\end{verbatim}
\subsubsection{\texttt{supported-value} — Value + Justification}
\label{sec:org3bcb6c4}

\begin{verbatim}
(struct supported-value
  (value      ;; the lattice value
   support)   ;; champ-root : assumption-id → #t (set of assumptions)
  #:transparent)
\end{verbatim}
\subsubsection{\texttt{tms-cell} — Truth-Maintained Cell (Immutable)}
\label{sec:org4c8c95c}

A TMS cell holds multiple contingent values, each justified by a different
assumption set. It generalizes a prop-cell to handle nondeterminism.

\begin{verbatim}
(struct tms-cell
  (values        ;; list of supported-value
   dependents)   ;; champ-root (set of prop-id → #t)
  #:transparent)

;; Note: immutable. Update returns new tms-cell via struct-copy.
;; The cell identity (tms-cell-id) is the key in the ATMS's cells map.
\end{verbatim}
\subsubsection{\texttt{atms} — The Persistent ATMS}
\label{sec:org246ed23}

\begin{verbatim}
(struct atms
  (network         ;; prop-network (the underlying persistent network)
   assumptions     ;; champ-root : assumption-id → assumption
   nogoods         ;; list of champ-root (each a set of assumption-ids → #t)
   tms-cells       ;; champ-root : cell-id → tms-cell
   next-assumption ;; Nat — monotonic counter
   believed)       ;; champ-root : assumption-id → #t (current worldview)
  #:transparent)
\end{verbatim}

Key design choices:
\begin{itemize}
\item \textbf{All CHAMP-backed}: assumptions, nogoods (as frozen sets), tms-cells, believed
\item \textbf{Worldview switching}: \texttt{(struct-copy atms a [believed new-set])} — O(1)
\item \textbf{Backtracking}: keep old \texttt{atms} reference — O(1), structural sharing
\item \textbf{No mutable state}: all operations return new \texttt{atms} values
\end{itemize}
\subsection{5.3 Pure Operations}
\label{sec:org3c9be2a}

\begin{verbatim}
;; Create a new assumption: returns (values new-atms assumption-id)
(define (atms-assume atms name datum)
  (define aid (assumption-id (atms-next-assumption atms)))
  (define a (assumption name datum))
  (values
   (struct-copy atms atms
     [assumptions (champ-insert (atms-assumptions atms)
                                 (equal-hash-code aid) aid a)]
     [next-assumption (+ 1 (atms-next-assumption atms))]
     [believed (champ-insert (atms-believed atms)
                              (equal-hash-code aid) aid #t)])
   aid))

;; Record a nogood: returns new atms with pruned worldview
(define (atms-add-nogood atms nogood-set)
  (struct-copy atms atms
    [nogoods (cons nogood-set (atms-nogoods atms))]))

;; Check consistency of an assumption set
(define (atms-consistent? atms assumption-set)
  (not (for/or ([ng (in-list (atms-nogoods atms))])
         (champ-subset? ng assumption-set))))

;; Switch worldview: returns new atms with different believed set
(define (atms-with-worldview atms new-believed)
  (struct-copy atms atms [believed new-believed]))
\end{verbatim}
\subsection{5.4 The \texttt{amb} Operator (Pure)}
\label{sec:org0c27f7f}

\texttt{amb} creates a choice point with n alternatives. Returns new ATMS.

\begin{verbatim}
(define (atms-amb atms alternatives)
  ;; 1. Create fresh hypotheses h1..hn
  (define-values (atms* hyps)
    (for/fold ([a atms] [hs '()])
              ([alt (in-list alternatives)]
               [i (in-naturals)])
      (define-values (a* hid) (atms-assume a (format "h~a" i) alt))
      (values a* (cons hid hs))))
  ;; 2. Record mutual exclusion: any pair is a nogood
  (define atms**
    (for*/fold ([a atms*])
               ([h1 (in-list hyps)]
                [h2 (in-list hyps)]
                #:when (not (equal? h1 h2)))
      (atms-add-nogood a (champ-insert (champ-insert champ-empty
                                          (equal-hash-code h1) h1 #t)
                                        (equal-hash-code h2) h2 #t))))
  ;; 3. Create TMS cell with supported values
  (values atms** hyps))
\end{verbatim}
\subsection{5.5 Two-Tier Mode: Lazy ATMS Activation}
\label{sec:orgcc5db7c}

The ATMS adds overhead for every cell operation: label management, support
set tracking, consistency checking. For \textbf{deterministic} queries (no \texttt{amb},
no choice points), this overhead is pure waste.

Design requirement: the ATMS operates in two tiers:

\textbf{Tier 1 — PropNetwork Only (default)}:
\begin{itemize}
\item No ATMS overhead. Just a \texttt{prop-network} with pure propagation.
\item Cells hold simple values, not supported values.
\item \texttt{run-to-quiescence} operates directly on the persistent network.
\item This is sufficient for Datalog evaluation, deterministic constraint
propagation, and simple tabled predicates.
\end{itemize}

\textbf{Tier 2 — Full ATMS (activated on first \texttt{amb})}:
\begin{itemize}
\item When the solver encounters its first choice point (\texttt{amb}), it upgrades
to full ATMS mode by wrapping the current network in an \texttt{atms} struct.
\item From that point on, all cell values become supported values (tagged with
assumption sets).
\item Nogoods, worldview switching, and DDB become available.
\end{itemize}

\begin{verbatim}
;; The solver state is either a bare network or a full ATMS
;; (struct solver-state (kind data) #:transparent)
;; kind is 'network | 'atms

;; On first amb:
(define (solver-upgrade-to-atms state)
  (if (eq? (solver-state-kind state) 'atms)
      state  ;; already upgraded
      (solver-state 'atms
        (atms-from-network (solver-state-data state)))))

;; atms-from-network : prop-network → atms
;; Wraps existing network, converting all cell values to
;; unconditionally supported values (support = ∅ = always true)
\end{verbatim}

The \texttt{solve-with} configuration supports:
\begin{itemize}
\item \texttt{:strategy :auto} (default) — starts in Tier 1, upgrades on first \texttt{amb}
\item \texttt{:strategy :depth-first} — pure PropNetwork, no ATMS, chronological choice
\item \texttt{:strategy :atms} — full ATMS from the start
\end{itemize}

This addresses the research doc's §14.1 concern: "For simple deterministic
programs, the ATMS layer is pure cost." Lazy activation means deterministic
programs never pay ATMS overhead.
\subsection{5.6 Contradiction Handler (Dependency-Directed Backtracking)}
\label{sec:org807f1aa}

When the underlying prop-network detects contradiction (a cell's merged
value = top):

\begin{enumerate}
\item Extract the support set from the contradicted cell's TMS values
\item Record as a nogood in the ATMS
\item The nogood automatically prunes that worldview from future exploration
\item Return new ATMS with nogood recorded
\end{enumerate}

This is pure: each step returns a new \texttt{atms} value. Dependency-directed
backtracking identifies \emph{which} choice was wrong, not just the most recent.
\subsection{5.7 Answer Collection (Pure)}
\label{sec:orgf134b3a}

\begin{verbatim}
;; Enumerate all consistent worldviews and extract answers
(define (atms-solve-all atms goal-cell-id)
  ;; 1. Enumerate maximal consistent assumption sets
  ;; 2. For each worldview, run the network to quiescence
  ;; 3. Extract goal cell value under that worldview
  ;; 4. Collect distinct answers
  ;; Returns: list of answer values
  ...)
\end{verbatim}
\subsection{5.8 AST Nodes (\textasciitilde{}10)}
\label{sec:org0c9c3b0}

\begin{center}
\begin{tabular}{lll}
Node & Fields & Semantics\\
\hline
\texttt{expr-atms-new} & network & Create ATMS wrapping network\\
\texttt{expr-amb} & atms, alternatives & Create choice point\\
\texttt{expr-assume} & atms, name, datum & Create hypothesis\\
\texttt{expr-retract} & atms, assumption-id & Retract hypothesis\\
\texttt{expr-nogood} & atms, assumption-set & Record nogood\\
\texttt{expr-solve-all} & atms, goal & Collect all solutions\\
\texttt{expr-tms-read} & atms, cell, worldview & Read under worldview\\
\texttt{expr-tms-write} & atms, cell, val, support & Write supported value\\
\texttt{expr-atms-type} & — & Type constructor ATMS\\
\texttt{expr-supported-type} & val-type & Type of SupportedValue\\
\end{tabular}
\end{center}
\subsection{5.9 New/Modified Files}
\label{sec:orgbfd5b88}

\begin{center}
\begin{tabular}{ll}
File & Changes\\
\hline
\texttt{racket/prologos/atms.rkt} & NEW: Persistent ATMS\\
\texttt{racket/prologos/syntax.rkt} & +10 AST nodes\\
\texttt{racket/prologos/propagator.rkt} & Contradiction → record in ATMS\\
+ standard pipeline files & Type rules, reduction, etc.\\
\texttt{tests/test-atms.rkt} & NEW\\
\texttt{tests/test-atms-search.rkt} & NEW\\
\texttt{tests/test-atms-backtrack.rkt} & NEW\\
\end{tabular}
\end{center}
\subsection{5.10 Tests (\textasciitilde{}50)}
\label{sec:org006a833}

\begin{itemize}
\item ATMS creation from network
\item Assumption creation: returns new atms + id
\item Supported value: value + support set
\item TMS cell: multiple contingent values
\item Consistency check: no nogood subsets
\item amb: creates mutual exclusion nogoods
\item amb: different alternatives under different worldviews
\item Nogood recording: prunes invalid worldviews
\item Dependency-directed backtracking: contradiction identifies cause
\item Worldview switching: \texttt{atms-with-worldview} — O(1)
\item Backtracking: keep old atms reference — O(1)
\item Persistence: old atms unchanged after assume/nogood
\item solve-all: enumerates all consistent answers
\item solve-all: deduplicates answers
\item Integration: persistent network + persistent ATMS
\item Performance: 100-alternative amb, solve-all
\end{itemize}
\subsection{5.11 Dependencies}
\label{sec:orgbd84cd5}

\begin{itemize}
\item Phase 2 (Persistent PropNetwork — ATMS wraps a network)
\item Phase 3 (PropNetwork as Prologos type — for typed ATMS operations)
\end{itemize}

---
\section{Phase 6: Tabling — SLG-Style Memoization}
\label{sec:org91408e1}

\subsection{6.1 Goal}
\label{sec:org3315cf2}

Implement tabling for completeness. Without tabling, left-recursive rules
cause infinite loops. Tabling memoizes intermediate results and detects
fixed-point completion.
\subsection{6.2 Design (XSB-Style SLG Resolution)}
\label{sec:org5f7aaea}

\begin{center}
\begin{tabular}{ll}
Concept & Implementation\\
\hline
Table & PropNetwork cell with SetLattice merge\\
Producer & Propagator computing new answers for table cell\\
Consumer & Propagator reading from table cell\\
Completion & Table cell quiescent when no new answers propagate\\
Answer mode & \texttt{all} (collect all) or \texttt{lattice} (join)\\
\end{tabular}
\end{center}
\subsection{6.3 Table Lifecycle}
\label{sec:orgf25ef77}

\begin{verbatim}
1. First call to tabled predicate → create table cell (SetLattice of answers)
2. Launch producer propagator → evaluates rules, writes answers to table cell
3. Recursive call to same predicate → create consumer propagator
   - Consumer reads from table cell
   - Re-fires when table cell's value grows (new answers appear)
4. Network reaches quiescence → no new answers propagate
5. Table cell value is the complete answer set
6. Subsequent calls → read from quiescent network (already computed)
\end{verbatim}
\subsection{6.4 Core Data Structures (Persistent)}
\label{sec:orgf3e4388}

The \texttt{table-store} is an \textbf{index}, not a store. It maps predicate call patterns
to cell-ids in the PropNetwork. The actual answer data lives in PropNetwork
cells with \texttt{SetLattice} merge — the table-store merely tells you \emph{which}
cell to look in.

\begin{verbatim}
;; In racket/prologos/tabling.rkt

(struct table-entry
  (predicate-name  ;; symbol
   call-pattern    ;; the specific instantiation pattern
   cell-id         ;; cell-id in the prop-network (answers live here)
   status)         ;; 'active | 'complete
  #:transparent)   ;; immutable — update via struct-copy

;; NOTE: no `answers` field — answers are stored in the PropNetwork cell
;; identified by `cell-id`. The cell uses SetLattice merge, so writing
;; a new answer to it computes set-union with existing answers.

(struct table-store
  (tables)     ;; champ-root : table-key → table-entry
  #:transparent)

;; Table key: predicate name + call pattern (up to renaming)
(define (table-key name args)
  (cons name (abstract-call-pattern args)))

;; Tables are backed by CHAMP for structural sharing.
;; Backtracking = keep old table-store reference + old prop-network.
;;
;; The lifecycle:
;; 1. First call → create table-entry + new cell in network
;; 2. Producer propagator writes answers to the cell
;; 3. Consumer propagators read from the cell (re-fire when it grows)
;; 4. Quiescence → mark table-entry status as 'complete
;; 5. Subsequent calls → read from the cell (already computed)
\end{verbatim}
\subsection{6.5 Lattice Answer Modes}
\label{sec:orga7a6b43}

Following XSB Prolog:

\begin{itemize}
\item \textbf{\texttt{all}}: Table stores set of all distinct answer substitutions
(\texttt{SetLattice} on substitutions)
\item \textbf{\texttt{lattice f}}: Table stores lattice join of all answers via \texttt{f}
(single aggregated value, new answers only "new" if they improve)
\item \textbf{\texttt{first}}: Table frozen after first answer (\texttt{once} semantics)
\end{itemize}
\subsection{6.6 Spec Metadata Integration}
\label{sec:org52aa3ba}

Tabling is declared via spec metadata:

\begin{verbatim}
spec ancestor : String -> String -> Prop
  :tabled true
  :answer-mode all        ;; default
\end{verbatim}

This requires adding \texttt{:tabled} and \texttt{:answer-mode} cases to
\texttt{parse-spec-metadata} in \texttt{macros.rkt} (following the \texttt{:examples} pattern).
\subsection{6.7 AST Nodes (\textasciitilde{}8)}
\label{sec:orga922177}

\begin{center}
\begin{tabular}{lll}
Node & Fields & Semantics\\
\hline
\texttt{expr-table-lookup} & name, args & Check if answer in table\\
\texttt{expr-table-new} & name, answer-mode & Create new table\\
\texttt{expr-table-add} & table, answer & Add answer to table\\
\texttt{expr-table-freeze} & table & Freeze table\\
\texttt{expr-table-answers} & table & Get all answers\\
\texttt{expr-table-complete?} & table & Check if table is complete\\
\texttt{expr-tabled} & spec-name & Mark predicate as tabled\\
\texttt{expr-table-type} & answer-type & Type of table\\
\end{tabular}
\end{center}
\subsection{6.8 New/Modified Files}
\label{sec:org62c1eb5}

\begin{center}
\begin{tabular}{ll}
File & Changes\\
\hline
\texttt{racket/prologos/tabling.rkt} & NEW: Table store, lifecycle\\
\texttt{racket/prologos/macros.rkt} & :tabled, :answer-mode metadata\\
\texttt{racket/prologos/syntax.rkt} & +8 AST nodes\\
+ standard pipeline files & Type rules, reduction, etc.\\
\texttt{tests/test-tabling.rkt} & NEW\\
\texttt{tests/test-tabling-lattice.rkt} & NEW\\
\end{tabular}
\end{center}
\subsection{6.9 Tests (\textasciitilde{}40)}
\label{sec:org0f9d435}

\begin{itemize}
\item Table creation: cell with SetLattice merge in PropNetwork
\item Table answer insertion: cell write grows answer set
\item Network quiescence: table complete when no new answers propagate
\item Producer: generates answers from rules, writes to table cell
\item Consumer: reads from table cell, re-fires on growth
\item Left-recursion: tabling terminates (the critical test)
\item Lattice answer mode: min aggregation
\item Lattice answer mode: set union aggregation
\item First answer mode: stops after one
\item Multiple tables: independent cells in same network
\item Backtracking: persistent network, old table state preserved
\item Spec metadata: :tabled parsed and stored
\item Integration: tabling + persistent ATMS for recursive search
\end{itemize}
\subsection{6.10 Dependencies}
\label{sec:org42e7ec7}

\begin{itemize}
\item Phase 3 (PropNetwork cells — tables are cells with SetLattice merge)
\item Phase 5 (ATMS — tabling + nondeterminism interact)
\end{itemize}

---
\section{Phase 7: Surface Syntax — \texttt{defr}, \texttt{rel}, \texttt{solve}, \texttt{\&>}}
\label{sec:orgddd09cd}

\subsection{7.1 Goal}
\label{sec:org3a0194a}

Implement the user-facing relational language as described in
\href{principles/RELATIONAL\_LANGUAGE\_VISION.org}{RELATIONAL\textsubscript{LANGUAGE}\textsubscript{VISION.org}}.
\subsection{7.2 Reader Changes}
\label{sec:org28c63a4}

The reader must handle:

\begin{center}
\begin{tabular}{ll}
Syntax & Reader Output\\
\hline
\texttt{?var} & \texttt{(logic-var var)}\\
\texttt{-var} & \texttt{(mode-var in var)}\\
\texttt{+var} & \texttt{(mode-var out var)}\\
\texttt{\&>} & \texttt{(\$clause-sep)}\\
\texttt{(goal ...)} & \texttt{(goal ...)} (parenthetical = relational)\\
\end{tabular}
\end{center}
\subsection{7.3 Parser Changes}
\label{sec:org2173161}

New surface AST nodes:

\begin{center}
\begin{tabular}{ll}
Form & Surface AST\\
\hline
\texttt{defr name [args] body} & \texttt{(surf-defr name args clauses)}\\
\texttt{(rel [args] body)} & \texttt{(surf-rel args clauses)}\\
\texttt{\&> g1 g2 ...} & \texttt{(surf-clause (g1 g2 ...))} within defr/rel\\
\texttt{(solve [goal])} & \texttt{(surf-solve goal)}\\
\texttt{(solve-with :opts goal)} & \texttt{(surf-solve-with opts goal)}\\
\texttt{(} ?x ?y)= & \texttt{(surf-unify x y)}\\
\texttt{(is ?x [expr])} & \texttt{(surf-is var expr)}\\
\end{tabular}
\end{center}
\subsection{7.4 Elaboration}
\label{sec:org7c55d2a}

Relations elaborate to propagator networks. Here is the concrete Racket-level
translation for both \texttt{defr} and \texttt{solve}:
\subsubsection{7.4.1 \texttt{defr} Elaboration}
\label{sec:orgaaa06e6}

\begin{verbatim}
;; defr ancestor [?x ?y]
;;   &> (parent ?x ?y)
;;   &> (parent ?x ?z) (ancestor ?z ?y)
;;
;; elaborates to:
;; 1. Create table for ancestor (tabled by default)
;; 2. For each clause, create a propagator that:
;;    a. Unifies arguments with clause head
;;    b. Resolves subgoals (recursively, via table)
;;    c. If all succeed, adds answer to table
;; 3. Register in relation store
\end{verbatim}
\subsubsection{7.4.2 \texttt{solve} Elaboration (Functional-Relational Bridge)}
\label{sec:orgd8e0a8a}

The \texttt{solve} form is the bridge from relational goals back to functional values.
Here is how it elaborates at the Racket level:

\begin{verbatim}
;; Source: functional code using solve
defn find-ancestors [person]
  (solve [ancestor person ?who])
;; Returns: Seq (Map Symbol Value)
\end{verbatim}

Elaborates to:

\begin{verbatim}
;; Racket-level translation of (solve [ancestor person ?who])
(define (solve-ancestor-query person)
  (let* (;; 1. Create solver state (network + uf-store)
         [net0 (make-prop-network)]
         [uf0  (uf-empty)]

         ;; 2. Create logic variable cells
         ;;    ?who is the query variable; person is ground (bound)
         [net1+who (net-new-cell net0 'bot flat-join flat-top?)]
         [net1 (car net1+who)]
         [who-id (cdr net1+who)]

         ;; 3. Look up 'ancestor in the relation registry
         ;;    (registered by defr elaboration)
         [ancestor-rel (relation-lookup 'ancestor)]

         ;; 4. Create query propagators:
         ;;    - Instantiate ancestor's clauses with (person, ?who)
         ;;    - Wire table cell for ancestor
         [net2 (relation-instantiate ancestor-rel net1 uf0
                  (list (ground-val person) who-id))]

         ;; 5. Run to quiescence
         [net3 (run-to-quiescence net2)]

         ;; 6. Extract answers from the table cell
         [answers (table-answers net3 'ancestor)]

         ;; 7. Project query variable bindings into Maps
         [results (for/list ([subst (in-set answers)])
                    (hash 'who (substitution-lookup subst who-id)))])
    ;; Return as lazy sequence of maps
    (list->lseq results)))
\end{verbatim}

Key points:
\begin{itemize}
\item \texttt{solve} creates a fresh \texttt{(network, uf-store)} pair scoped to the query
\item Ground arguments are injected directly; logic variables become cells
\item The relation's propagators run to quiescence in the fresh network
\item Results are projected from table cells into \texttt{Seq (Map Symbol Value)}
\item The entire operation is \textbf{pure} — no side effects, no global state
\end{itemize}
\subsection{7.5 Grammar Updates}
\label{sec:org6482395}

Both grammar files must be updated:
\begin{itemize}
\item \texttt{docs/spec/grammar.ebnf} — EBNF production rules
\item \texttt{docs/spec/grammar.org} — Prose companion with examples
\end{itemize}

New productions:

\begin{verbatim}
relation-def  = "defr" , identifier , param-list , clause+ ;
anonymous-rel = "(" , "rel" , param-list , clause+ , ")" ;
clause        = "&>" , goal+ ;
goal          = "(" , goal-head , goal-arg* , ")" ;
goal-head     = identifier | "=" | "is" | "not" ;
goal-arg      = logic-var | expression ;
logic-var     = "?" , identifier ;
mode-var      = ( "-" | "+" ) , identifier ;
solve-expr    = "(" , "solve" , "[" , goal , "]" , ")" ;
solve-with    = "(" , "solve-with" , metadata* , "[" , goal , "]" , ")" ;
\end{verbatim}
\subsection{7.6 Compile-Time Stratification Check}
\label{sec:org8aa8e0f}

Any program using \texttt{not} (negation-as-failure) must pass a stratification check
at compile time. This is the minimum viable safety guarantee for negation:

\begin{enumerate}
\item \textbf{Build the predicate dependency graph}: For each \texttt{defr}, record which
predicates it calls positively and negatively (via \texttt{not}).

\item \textbf{Compute SCCs}: Strongly connected components via Tarjan's algorithm.
Each SCC is a group of mutually recursive predicates.

\item \textbf{Classify negative edges}: An edge \texttt{A -{}-{}not-{}-{}> B} is:
\begin{itemize}
\item \textbf{Safe} if \texttt{B} is in a \emph{lower} SCC than \texttt{A} (B is fully computed before
A uses its negation)
\item \textbf{Unsafe} if \texttt{B} is in the \emph{same} SCC as \texttt{A} (cyclic negation — the
truth of \texttt{not B} depends on \texttt{A} which depends on \texttt{not B})
\end{itemize}

\item \textbf{Reject unstratifiable programs}: If any negative edge is within an SCC,
emit a compile-time error: "Cyclic negation detected: \texttt{A} and \texttt{B} are
mutually recursive through negation. Consider restructuring."

\item \textbf{Assign strata}: Each SCC gets a stratum number (topological order of
the condensation DAG). Lower strata are evaluated first; their results
are frozen before higher strata use them via \texttt{not}.
\end{enumerate}

\begin{verbatim}
;; In racket/prologos/stratify.rkt (NEW)

;; stratify : (Listof defr-info) → (Listof (Listof defr-info))
;; Returns a list of strata (each stratum = list of defrs to evaluate together)
;; Raises compile-time error if program is unstratifiable

(define (stratify defr-infos)
  (define dep-graph (build-dependency-graph defr-infos))
  (define sccs (tarjan-scc dep-graph))
  (define condensation (condense dep-graph sccs))
  ;; Check: no negative edge within any SCC
  (for ([scc (in-list sccs)])
    (for ([pred (in-list scc)])
      (for ([neg-dep (in-list (negative-dependencies pred))])
        (when (member neg-dep scc)
          (raise-stratification-error pred neg-dep scc)))))
  ;; Return strata in topological order
  (topological-sort condensation))
\end{verbatim}

This check runs during elaboration of \texttt{defr} forms. Programs without \texttt{not} are
trivially stratifiable (single stratum) and incur no analysis cost.
\subsection{7.7 AST Nodes (\textasciitilde{}15)}
\label{sec:org3c1b050}

\begin{center}
\begin{tabular}{lll}
Node & Fields & Semantics\\
\hline
\texttt{expr-defr} & name, params, clauses & Named relation\\
\texttt{expr-rel} & params, clauses & Anonymous relation\\
\texttt{expr-clause} & goals & Single clause (disjunct)\\
\texttt{expr-goal-app} & name, args & Relational goal application\\
\texttt{expr-logic-var} & name, mode & Logic variable\\
\texttt{expr-unify-goal} & lhs, rhs & Unification goal (= ?x ?y)\\
\texttt{expr-is-goal} & var, expr & Functional evaluation in relation\\
\texttt{expr-not-goal} & goal & Negation-as-failure\\
\texttt{expr-solve} & goal & Solve returning Seq Map\\
\texttt{expr-solve-with} & opts, goal & Parameterized solve\\
\texttt{expr-solve-one} & goal & Solve returning first answer\\
\texttt{expr-relation-type} & param-types & Type of a relation\\
\texttt{expr-goal-type} & — & Type of a goal\\
\texttt{expr-cut} & — & Committed choice (once)\\
\texttt{expr-guard} & condition, goal & Guard evaluation\\
\end{tabular}
\end{center}
\subsection{7.8 New/Modified Files}
\label{sec:org45714a4}

\begin{center}
\begin{tabular}{ll}
File & Changes\\
\hline
\texttt{racket/prologos/reader.rkt} & ?var, -var, +var, \&> handling\\
\texttt{racket/prologos/surface-syntax.rkt} & New surface AST structs\\
\texttt{racket/prologos/parser.rkt} & Parse defr, rel, solve, \&>\\
\texttt{racket/prologos/macros.rkt} & process-defr, process-rel\\
\texttt{racket/prologos/elaborator.rkt} & Elaborate relations → propagators\\
\texttt{racket/prologos/stratify.rkt} & NEW: SCC + stratification check\\
\texttt{racket/prologos/syntax.rkt} & +15 AST nodes\\
\texttt{racket/prologos/typing-core.rkt} & Type rules for relational forms\\
+ standard pipeline files & All 12+ consuming modules\\
\texttt{docs/spec/grammar.ebnf} & New productions\\
\texttt{docs/spec/grammar.org} & New sections\\
\texttt{lib/prologos/core/relations.prologos} & NEW: Standard relations\\
\texttt{tests/test-relations-basic.rkt} & NEW\\
\texttt{tests/test-relations-tabling.rkt} & NEW\\
\texttt{tests/test-relations-search.rkt} & NEW\\
\texttt{tests/test-relations-ws.rkt} & NEW: WS-mode integration\\
\texttt{tests/test-solve.rkt} & NEW\\
\end{tabular}
\end{center}
\subsection{7.9 Tests (\textasciitilde{}80)}
\label{sec:org22b49a3}

\begin{itemize}
\item Reader: ?var, -var, +var parsed correctly
\item Parser: defr, rel, \&>, solve parsed correctly
\item WS-mode: defr with indentation-based clauses
\item Sexp-mode: defr with explicit (defr \ldots{}) form
\item Basic relation: parent facts, query
\item Recursive relation: ancestor with tabling
\item Unification goal: = ?x ?y
\item Functional evaluation: is ?x [expr]
\item Multiple clauses: \&> separator
\item Anonymous relation: (rel [?x ?y] \&> \ldots{})
\item Solve: returns Seq of substitution maps
\item Solve: empty result for unsatisfiable goal
\item Solve-with: :strategy option
\item Solve-with: :timeout option
\item Negation: not (goal)
\item Mode annotations: -var, +var optimization hints
\item Integration with functional code: solve in defn body
\item Integration with traits: relation using trait methods
\item Stratification: non-negated program = single stratum (trivial)
\item Stratification: negated program with valid stratification compiles
\item Stratification: cyclic negation detected and rejected at compile time
\item Error messages: undefined relation, arity mismatch
\item Performance: 100-fact database, recursive query
\end{itemize}
\subsection{7.10 Dependencies}
\label{sec:orgc50449a}

\begin{itemize}
\item ALL previous phases (Lattice, Cells, LVars, UF, ATMS, Tabling)
\end{itemize}

---
\section{Phase Summary}
\label{sec:org6675575}

\begin{center}
\begin{tabular}{rlrrrll}
Phase & Name & New AST & New Files & New Tests & Deps & Size\\
\hline
1 & Lattice + champ-insert-join & 0 & 3 & 25 & None & Small\\
2 & Persistent PropNetwork & 0 & 3 & 60 & Ph 1 & Medium\\
3 & PropNetwork Prologos Type & 12 & 3 & 50 & Ph 2 & Medium\\
4 & UnionFind & 6 & 2 & 30 & None (∥3) & Small\\
5 & Persistent ATMS & 10 & 3 & 50 & Ph 3 & Large\\
6 & Tabling & 8 & 3 & 40 & Ph 3,5 & Medium\\
7 & Surface Syntax & 15 & 11+ & 83 & All & Large\\
\hline
TOTAL &  & 51 & 28+ & 338 &  & \\
\end{tabular}
\end{center}

Estimated total: \textasciitilde{}51 new AST nodes (was 61 — LVar nodes eliminated),
\textasciitilde{}28+ new files, \textasciitilde{}338 new tests.

Key changes from the original design:
\begin{itemize}
\item \textbf{Persistent architecture}: All structs are \texttt{\#:transparent} (not \texttt{\#:mutable})
\item \textbf{LVars subsumed}: 10 LVar AST nodes eliminated (LVars = network cells)
\item \textbf{Phase 2 split}: Racket-level implementation (Phase 2) separated from
Prologos type exposure (Phase 3)
\item \textbf{Free backtracking}: O(1) via structural sharing (CHAMP maps)
\item \textbf{Free snapshots}: Networks are values — keeping a reference IS a snapshot
\item \textbf{Per-cell contradiction detection}: \texttt{contradicts?} predicate per cell, not
in Lattice trait (no trait hierarchies — see DESIGN\textsubscript{PRINCIPLES.org})
\item \textbf{ATMS two-tier mode}: Lazy activation on first \texttt{amb} — deterministic
programs never pay ATMS overhead
\item \textbf{Compile-time stratification}: SCC + negative edge check rejects
unstratifiable programs at compile time
\end{itemize}

---
\section{Interaction with Existing Infrastructure}
\label{sec:org11fd55e}

\subsection{Metavar System}
\label{sec:orgc312453}

The current metavar system (\texttt{current-meta-store}, \texttt{save/restore-meta-state!})
is \emph{not} replaced by Phase 2's propagator network. The elaborator continues
to use its existing metavar system for type inference. The propagator network
is a \emph{separate, parallel} system used by the logic engine at runtime.

The persistent PropNetwork architecture was directly motivated by the
\texttt{save/restore-meta-state!} problem: the metavar store requires O(n) deep
copies for speculative type checking, and restore requires explicit undo.
The persistent approach avoids this entirely — backtracking is O(1) (keep
old reference). In the future (post-Phase 7), the elaborator's metavar
system could be refactored to use propagator cells internally.
\subsection{Trait System}
\label{sec:org4318422}

The \texttt{Lattice} trait (Phase 1) uses the existing trait infrastructure with
no modifications. Lattice instances resolve via \texttt{resolve-trait-constraints!}
like any other trait.
\subsection{Spec Metadata}
\label{sec:org65218bc}

Phase 6 adds \texttt{:tabled} and \texttt{:answer-mode} to \texttt{parse-spec-metadata}. This
follows the same pattern as \texttt{:examples} (Stage C of Extended Spec Hardening):
explicit case in the metadata parser using \texttt{collect-constraint-values} or
direct value capture.
\subsection{Warnings}
\label{sec:org77d1bd3}

The logic engine may emit new warning types (e.g., "tabled predicate exceeded
table size limit", "negation in unstratifiable position"). These follow the
\texttt{warnings.rkt} pattern: new struct + parameter + emit/format.
\subsection{QTT / Multiplicities}
\label{sec:org7f0bfa5}

Logic variables live at multiplicity \texttt{:w} (unrestricted). They are shared
across multiple goals and clauses. The binding environment (the substitution
store) is functional/persistent and does not consume resources linearly.

Open question: should cells be linear (\texttt{:1})? A cell created and consumed
exactly once (write, then read, then discard) could be linear. For now,
cells are unrestricted (\texttt{:w}).
\subsection{Collections}
\label{sec:org20adc7e}

The entire propagator network is backed by CHAMP maps (\texttt{champ.rkt}):
\begin{itemize}
\item \texttt{cells} map: \texttt{cell-id → prop-cell}
\item \texttt{propagators} map: \texttt{prop-id → propagator}
\item \texttt{merge-fns} map: \texttt{cell-id → merge-fn}
\item ATMS \texttt{assumptions}, \texttt{tms-cells}, \texttt{believed}: all CHAMP maps
\item Table stores: CHAMP map of table entries
\end{itemize}

The \texttt{champ-insert-join} helper (added in Phase 1) enables lattice-aware
insert-with-merge at the CHAMP level. SetLattice cells use \texttt{set-union}
as their merge function. MapLattice cells use \texttt{merge-with join}.

No separate LVar-Set or LVar-Map types are needed — these are simply
cells in a PropNetwork with appropriate merge functions.

---
\section{Appendix A: Resolution by Example — \texttt{ancestor} as Propagators}
\label{sec:orgabcb971}

This appendix shows the complete elaboration of a \texttt{defr} definition into
propagator network operations. It bridges the gap between the surface syntax
(Phase 7) and the propagator substrate (Phases 2-6).
\subsection{A.1 Source Program}
\label{sec:org5107f9b}

\begin{verbatim}
;; Facts
defr parent [?x ?y]
  &> (= ?x "alice") (= ?y "bob")
  &> (= ?x "bob") (= ?y "carol")
  &> (= ?x "bob") (= ?y "dave")

;; Recursive relation (tabled by default)
defr ancestor [?x ?y]
  &> (parent ?x ?y)                      ;; clause 1: base case
  &> (parent ?x ?z) (ancestor ?z ?y)    ;; clause 2: recursive
\end{verbatim}
\subsection{A.2 Step 1: Table Creation}
\label{sec:org2af97ad}

\texttt{ancestor} is tabled by default. The elaborator:

\begin{enumerate}
\item Creates a table entry in the \texttt{table-store} index:
\texttt{(table-entry 'ancestor <call-pattern> cell-42 'active)}
\item Creates a PropNetwork cell for the table's answer set:
\texttt{(net-new-cell net 'empty-set set-union)}  → \texttt{cell-42}
\item The cell's merge function is \texttt{set-union} (\texttt{SetLattice}) — answers accumulate.
\end{enumerate}
\subsection{A.3 Step 2: Clause 1 → Producer Propagator}
\label{sec:orgcaff9f1}

Clause 1: \texttt{\&> (parent ?x ?y)} (base case)

The elaborator creates a propagator that:

\begin{verbatim}
;; Clause-1 fire-fn for ancestor(?x, ?y):
;;   For each parent fact matching (?x, ?y), write to ancestor's table cell
(define (make-ancestor-clause1-fire-fn table-cell-id)
  (lambda (net)
    ;; For each fact in parent's table:
    (define parent-answers
      (net-cell-read net parent-table-cell-id))
    ;; Each parent answer is an ancestor answer (base case)
    ;; Write each to the ancestor table cell
    (for/fold ([net* net])
              ([ans (in-set parent-answers)])
      (net-cell-write net* table-cell-id
                      (set-add (set) ans)))))
\end{verbatim}

Inputs: \texttt{parent}'s table cell.
Outputs: \texttt{ancestor}'s table cell.
\subsection{A.4 Step 3: Clause 2 → Producer Propagator}
\label{sec:org89320c7}

Clause 2: \texttt{\&> (parent ?x ?z) (ancestor ?z ?y)} (recursive case)

This creates a propagator that:

\begin{verbatim}
;; Clause-2 fire-fn for ancestor(?x, ?y):
;;   Join parent(?x, ?z) with ancestor(?z, ?y) → ancestor(?x, ?y)
(define (make-ancestor-clause2-fire-fn table-cell-id)
  (lambda (net)
    (define parent-facts (net-cell-read net parent-table-cell-id))
    (define ancestor-facts (net-cell-read net table-cell-id))
    ;; For each parent(x, z) and ancestor(z, y), derive ancestor(x, y)
    (define new-answers
      (for*/set ([p (in-set parent-facts)]
                 [a (in-set ancestor-facts)]
                 #:when (equal? (subst-lookup p '?y)    ;; z in parent
                                (subst-lookup a '?x)))  ;; z in ancestor
        (make-subst '?x (subst-lookup p '?x)
                    '?y (subst-lookup a '?y))))
    (net-cell-write net table-cell-id new-answers)))
\end{verbatim}

Inputs: \texttt{parent}'s table cell AND \texttt{ancestor}'s table cell (self-reference!).
Outputs: \texttt{ancestor}'s table cell.

The self-reference is what makes tabling essential: the propagator reads from
the cell it writes to. Tabling ensures this reaches a fixed point rather
than looping infinitely.
\subsection{A.5 Step 4: Run to Quiescence}
\label{sec:org506b0a0}

\begin{verbatim}
Iteration 1:
  parent table cell = {(alice,bob), (bob,carol), (bob,dave)}
  clause-1 fires → ancestor cell ∪= {(alice,bob), (bob,carol), (bob,dave)}
  clause-2 fires → join parent × ancestor:
    parent(alice,bob) ∧ ancestor(bob,carol) → ancestor(alice,carol)
    parent(alice,bob) ∧ ancestor(bob,dave) → ancestor(alice,dave)
  ancestor cell ∪= {(alice,carol), (alice,dave)}

Iteration 2:
  clause-2 fires again (ancestor cell grew):
    parent(alice,bob) ∧ ancestor(bob,carol) → ancestor(alice,carol) — already present
    parent(alice,bob) ∧ ancestor(bob,dave) → ancestor(alice,dave) — already present
    No new answers added.

Quiescence reached. ancestor cell = {(alice,bob), (bob,carol), (bob,dave),
                                      (alice,carol), (alice,dave)}
\end{verbatim}
\subsection{A.6 Step 5: Answer Extraction}
\label{sec:orgd5cf1c9}

A query \texttt{(solve [ancestor "alice" ?who])} reads the ancestor table cell,
filters for entries where \texttt{?x = "alice"}, and projects the \texttt{?y} values:

\begin{verbatim}
Results: ?who ∈ {"bob", "carol", "dave"}
\end{verbatim}

---
\section{Appendix B: End-to-End Query Walkthrough (All Three Layers)}
\label{sec:orgb80a646}

This appendix shows a single query that exercises all three layers of the
architecture: PropNetwork (Layer 1), ATMS (Layer 2), and Stratification
(Layer 3).
\subsection{B.1 Source Program}
\label{sec:orgd43a7a1}

\begin{verbatim}
defr edge [?from ?to]
  &> (= ?from "a") (= ?to "b")
  &> (= ?from "b") (= ?to "c")
  &> (= ?from "a") (= ?to "c")

defr reachable [?x ?y]
  &> (edge ?x ?y)
  &> (edge ?x ?z) (reachable ?z ?y)

;; Negation triggers stratification (Layer 3)
defr unreachable [?x ?y]
  &> (node ?x) (node ?y) (not (reachable ?x ?y))

;; Query with nondeterminism (triggers ATMS, Layer 2)
;; "Find a node that is unreachable from 'a'"
let result := (solve [unreachable "a" ?target])
\end{verbatim}
\subsection{B.2 Compile-Time: Stratification Check}
\label{sec:org0514c30}

The predicate dependency graph:

\begin{verbatim}
edge         (no deps)         → Stratum 0
reachable    → edge (+)        → Stratum 0 (same SCC, all positive)
node         (no deps)         → Stratum 0
unreachable  → reachable (−)   → Stratum 1 (negative edge crosses strata)
\end{verbatim}

The stratification checker (§7.6) verifies: the negative edge \texttt{unreachable -{}-{}not-{}-{}> reachable} crosses from Stratum 1 to Stratum 0 — safe.
\subsection{B.3 Runtime: Layer 1 — PropNetwork (Stratum 0)}
\label{sec:orgfc35f7a}

Evaluate \texttt{edge}, \texttt{reachable}, and \texttt{node} to fixed point:

\begin{verbatim}
1. edge table cell:
   {(a,b), (b,c), (a,c)}

2. reachable table cell (iterate to fixpoint):
   Iteration 1: {(a,b), (b,c), (a,c)} — from edge base case
   Iteration 2: {(a,b), (b,c), (a,c), (a,c)} — (a,b)+(b,c)→(a,c), already present
   Fixpoint: {(a,b), (b,c), (a,c)}

3. node table cell (derived from edge endpoints):
   {a, b, c}
\end{verbatim}

All purely monotonic. No ATMS overhead. Layer 1 only.
\subsection{B.4 Runtime: Layer 3 — Stratum Boundary}
\label{sec:org08c45ba}

Stratum 0 is complete. Now the runtime evaluates Stratum 1.

\texttt{unreachable} uses \texttt{not (reachable ?x ?y)}. This reads the \emph{frozen} reachable
table from Stratum 0:

\begin{verbatim}
For each pair (x, y) in node × node:
  (a, a): reachable(a, a) = not found → unreachable(a, a) ✓
  (a, b): reachable(a, b) = found → skip
  (a, c): reachable(a, c) = found → skip
  (b, a): reachable(b, a) = not found → unreachable(b, a) ✓
  (b, b): reachable(b, b) = not found → unreachable(b, b) ✓
  (b, c): reachable(b, c) = found → skip
  (c, a): reachable(c, a) = not found → unreachable(c, a) ✓
  (c, b): reachable(c, b) = not found → unreachable(c, b) ✓
  (c, c): reachable(c, c) = not found → unreachable(c, c) ✓
\end{verbatim}
\subsection{B.5 Runtime: Layer 2 — ATMS (if needed)}
\label{sec:orga25bb78}

In this example, the query is deterministic (no \texttt{amb}). Under the two-tier
mode (§5.5), the solver stays in Tier 1 (PropNetwork only) — no ATMS
overhead.

If we modified the program to include a choice point:

\begin{verbatim}
;; "Pick a starting node and find what's unreachable from it"
defr mystery-query [?start ?target]
  &> (node ?start) (unreachable ?start ?target)
\end{verbatim}

Now \texttt{node ?start} produces 3 alternatives → \texttt{amb} activates, upgrading
to Tier 2 (full ATMS):

\begin{verbatim}
Hypothesis h1: ?start = "a"
  → unreachable("a", ?target) under h1
  → Answers: {(a,a)} under h1

Hypothesis h2: ?start = "b"
  → unreachable("b", ?target) under h2
  → Answers: {(b,a), (b,b)} under h2

Hypothesis h3: ?start = "c"
  → unreachable("c", ?target) under h3
  → Answers: {(c,a), (c,b), (c,c)} under h3

All hypotheses are consistent (no nogoods).
Final answers: union of all worldviews.
\end{verbatim}
\subsection{B.6 Summary: Which Layer Handles What}
\label{sec:org39c275f}

\begin{center}
\begin{tabular}{lll}
Query Aspect & Layer Used & Phase\\
\hline
\texttt{edge}, \texttt{reachable} tables & 1: PropNetwork & Ph 2,6\\
\texttt{not (reachable ...)} & 3: Stratification & Ph 7\\
Stratum evaluation order & 3: Stratification & Ph 7\\
\texttt{node ?start} choice points & 2: ATMS & Ph 5\\
Backtracking on contradictions & 2: ATMS & Ph 5\\
Table answer accumulation & 1: PropNetwork & Ph 2,6\\
\end{tabular}
\end{center}

---
\section{Performance Expectations}
\label{sec:org28b7347}

\subsection{Cost Model}
\label{sec:org453fa7f}

The persistent architecture adds a constant factor vs. mutable implementations.
Here are the expected costs for each operation:

\begin{center}
\begin{tabular}{llll}
Operation & Cost (Persistent) & Cost (Mutable) & Ratio\\
\hline
Cell read & O(log₃₂ n) ≈ O(C), C≤7 & O(1) hash lookup & \textasciitilde{}3x\\
Cell write (no change) & O(log₃₂ n) comparison & O(1) & \textasciitilde{}3x\\
Cell write (changed) & O(log₃₂ n) + alloc path & O(1) & \textasciitilde{}5x\\
Add propagator & O(log₃₂ n) × inputs & O(1) × inputs & \textasciitilde{}3-5x\\
Backtrack & O(1) — keep old ref & O(n) — deep copy & 100x+ better\\
Snapshot & O(1) — identity & O(n) — deep copy & 100x+ better\\
run-to-quiescence & O(k × log₃₂ n) per step & O(k) per step & \textasciitilde{}3-5x\\
\end{tabular}
\end{center}

Where:
\begin{itemize}
\item n = number of cells in the network (typically 100-10,000)
\item k = number of propagator firings to quiescence
\item log₃₂(10,000) ≈ 2.7 — so the "constant" is truly small
\end{itemize}
\subsection{Benchmark Targets (to validate during implementation)}
\label{sec:org2a5ff61}

\begin{center}
\begin{tabular}{lll}
Benchmark & Target & Validates\\
\hline
Create 1000-cell network & < 50ms & CHAMP scaling\\
Write 1000 values, run to quiescence & < 200ms & Propagation\\
Backtrack 100 times over 1000-cell net & < 1ms total & Persistence\\
100-fact ancestor query (tabled) & < 100ms & Tabling\\
1000-fact Datalog transitive closure & < 500ms & Scalability\\
\texttt{solve} with 100 alternatives (ATMS) & < 1s & ATMS overhead\\
Same 100 alternatives, Tier 1 only & < 200ms & Two-tier skip\\
\end{tabular}
\end{center}

These targets are order-of-magnitude estimates. Actual benchmarks will be
added to \texttt{tools/benchmark-tests.rkt} as each phase is implemented. The
key comparison is: persistent overhead vs. backtracking savings. For any
program with search (>1 choice point), the persistent approach wins
because each backtrack is O(1) instead of O(n) deep copy.
\subsection{When Performance Matters Most}
\label{sec:org5c76739}

For the logic engine's intended use cases:
\begin{itemize}
\item \textbf{Type checking}: The elaborator's current metavar system is fast for
deterministic type inference. The PropNetwork adds overhead here. The
elaborator refactoring (post-Phase 7) should only proceed after
benchmarking confirms acceptable overhead.
\item \textbf{Logic queries}: Search-heavy programs (Prolog-style) benefit enormously
from free backtracking. The persistent overhead is amortized over the
many backtracks avoided.
\item \textbf{Datalog}: Bottom-up evaluation with tabling is inherently monotonic —
no backtracking needed. Performance depends on propagation efficiency
and set operations.
\end{itemize}

---
\section{What This Design Does NOT Cover}
\label{sec:org0bf976f}

\subsection{Elaborator Refactoring (Phase 1 of Research Doc)}
\label{sec:orgb235e87}

Refactoring \texttt{current-meta-store} to use a persistent PropNetwork internally
is recommended by the research document but deferred. The persistent network
architecture makes this refactoring simpler: replace the mutable meta store
with a PropNetwork value threaded through elaboration. But it is still a
large internal change that can proceed independently after the logic engine
exists.
\subsection{Galois Connections / Domain Embeddings (Phase 6 of Research Doc)}
\label{sec:org8cfebf5}

Modular constraint domains connected via Galois connections are an advanced
feature deferred until the basic engine is operational.
\subsection{Full Stratified Evaluation Runtime (Phase 4 of Research Doc)}
\label{sec:orgd8a94ed}

Phase 7 includes a compile-time stratification \emph{check} (§7.6): SCC decomposition,
negative edge classification, and rejection of unstratifiable programs. The
runtime evaluation of strata (evaluating lower strata to fixpoint before
proceeding to higher strata) is also included as part of the solver's query
evaluation loop. Full support for aggregation operators (\texttt{count}, \texttt{min}, \texttt{max},
\texttt{sum}) between strata is deferred to a follow-up phase.
\subsection{CRDTs / Distributed Logic}
\label{sec:orgb08223d}

CRDT-backed collections for distributed actors are long-term goals not
addressed here.
\subsection{QuickCheck / Property Testing}
\label{sec:org586f1d0}

Executing \texttt{:holds} clauses and \texttt{:examples} entries requires the logic engine
(for proof search). Once the engine exists, this becomes a natural Phase 2
of the Extended Spec Design.

---
\section{Architectural Decision: Persistent Networks}
\label{sec:org51784be}

\subsection{The Problem with Mutable Propagator Networks}
\label{sec:org4e63d57}

The original design proposed mutable Racket structs for PropCell, PropNetwork,
TMSCell, and ATMS. This creates the same save/restore problem that plagues
the current metavar store:

\begin{itemize}
\item \textbf{Snapshotting} requires deep copies — O(n) for n cells
\item \textbf{Backtracking} requires explicit undo — O(mutations) to restore
\item \textbf{Network mobility} requires serialization — custom ser/deser code
\item \textbf{Debugging} is harder — state is invisible without explicit inspection
\end{itemize}
\subsection{The Persistent Solution}
\label{sec:org3304796}

Making the entire network a persistent value backed by CHAMP maps:

\begin{itemize}
\item \textbf{Snapshotting} = free (keep a reference — structural sharing)
\item \textbf{Backtracking} = O(1) (use old reference)
\item \textbf{Network mobility} = serialize the value (all data is in the struct)
\item \textbf{Debugging} = the network IS the state (print it, compare it)
\item \textbf{Cost}: O(log₃₂ n) per cell write instead of O(1)
\end{itemize}

Since log₃₂(n) ≤ 7 for any practical n (up to \textasciitilde{}34 billion cells),
the cost is bounded by a small constant. For the typical logic engine
workload (hundreds to thousands of cells), this is 2-3 levels of
CHAMP trie traversal — effectively O(1).
\subsection{LVar Elimination}
\label{sec:org7eb2fd6}

A further simplification: LVars are subsumed by PropNetwork cells.
An LVar is just a cell-id in a persistent prop-network, since
\texttt{net-cell-write} already does join-on-write (inflationary) via the
per-cell \texttt{merge-fn}. This eliminates \textasciitilde{}10 dedicated LVar AST nodes
and the separate \texttt{lvar.rkt} module.

\begin{center}
\begin{tabular}{ll}
Before & After\\
\hline
61 AST & 51 AST\\
32 files & 27+ files\\
Separate LVar, LVar-Set, LVar-Map & Cells with SetLattice, MapLattice\\
Mutable cells + scheduler & Persistent network, pure loop\\
\texttt{save/restore} for backtracking & Keep old reference\\
\end{tabular}
\end{center}

---
\section{Key Lessons from Prior Work}
\label{sec:orgf2a0100}

(Gathered from CLAUDE.md MEMORY and implementation experience)

\begin{enumerate}
\item \textbf{12+ file pipeline for new AST nodes} — every new \texttt{expr-*} touches
syntax.rkt, typing-core.rkt, reduction.rkt, elaborator.rkt, zonk.rkt,
pretty-print.rkt, qtt.rkt, substitution.rkt, unify.rkt, + tests.

\item \textbf{Per-command parameter isolation} — all runtime state (meta stores,
warning accumulators, table stores) must be parameterized and reset
per-command in \texttt{driver.rkt}.

\item \textbf{Trait dictionary dispatch} — single-method traits: dict IS the function.
Multi-method: nested Sigma. Dict names: \texttt{TypeArg-{}-{}TraitName-{}-{}dict}.

\item \textbf{\texttt{collect-constraint-values}} — the right collector for list-valued
metadata like \texttt{:examples}, \texttt{:laws}, \texttt{:tabled}.

\item \textbf{Generic \texttt{rewrite-implicit-map} branch} — WS-mode keyword blocks need
to call \texttt{rewrite-implicit-map} at dispatch. The property/functor bugs
taught this lesson.

\item \textbf{Transient pattern} — generic wrappers (\texttt{expr-transient}, \texttt{expr-persist})
\begin{itemize}
\item type-specific variants. Follow this for cell/lvar generics.
\end{itemize}

\item \textbf{Test file size} — keep under \textasciitilde{}20 cases / \textasciitilde{}30s for parallel execution.
Split large test files into -01, -02 parts.

\item \textbf{Update dep-graph.rkt} — every new .rkt file needs an entry in
\texttt{tools/dep-graph.rkt} for incremental testing to work.
\end{enumerate}

---
\section{References}
\label{sec:org4b88124}

Organized by phase relevance:
\subsection{Phase 1 (Lattice)}
\label{sec:orgbdb3d14}
\begin{itemize}
\item Cousot \& Cousot, "Abstract Interpretation: A Unified Lattice Model" (POPL, 1977)
\item Tarski, "A Lattice-Theoretical Fixpoint Theorem" (Pacific J. Math., 1955)
\end{itemize}
\subsection{Phase 2 (Propagators)}
\label{sec:orgf08fc2c}
\begin{itemize}
\item Radul \& Sussman, "The Art of the Propagator" (MIT TR, 2009)
\item Radul, \emph{Propagation Networks} (PhD thesis, MIT, 2009)
\item Hellerstein, "Keeping CALM" (CACM, 2020)
\end{itemize}
\subsection{Phase 3 (LVars)}
\label{sec:org8937437}
\begin{itemize}
\item Kuper \& Newton, "LVars: Lattice-Based Data Structures for Deterministic Parallelism" (FHPC, 2013)
\item Kuper et al., "Freeze After Writing" (POPL, 2014)
\end{itemize}
\subsection{Phase 4 (UnionFind)}
\label{sec:orgbbf0423}
\begin{itemize}
\item Conchon \& Filliâtre, "A Persistent Union-Find Data Structure" (ML Workshop, 2007)
\end{itemize}
\subsection{Phase 5 (ATMS)}
\label{sec:orgb61ef44}
\begin{itemize}
\item de Kleer, "An Assumption-Based TMS" (AI Journal, 1986)
\end{itemize}
\subsection{Phase 6 (Tabling)}
\label{sec:org55c3272}
\begin{itemize}
\item Swift \& Warren, "XSB: Extending Prolog with Tabled Logic Programming" (TPLP, 2012)
\item Chen \& Warren, "Tabled Evaluation with Delaying" (JACM, 1996)
\item Madsen et al., "From Datalog to Flix" (PLDI, 2016)
\end{itemize}
\subsection{Phase 7 (Surface)}
\label{sec:org4ff188f}
\begin{itemize}
\item Arntzenius \& Krishnaswami, "Datafun: A Functional Datalog" (ICFP, 2016)
\item Fruhwirth, \emph{Constraint Handling Rules} (CUP, 2009)
\end{itemize}
\end{document}
